 \paragraph*{04.05.08.01 Strategische Ressourcenplanung entwickeln, um die Projektergebnisse liefern zu
können}
\kcilabel{para:04.05.08.01_StrategischeRessourcenplanung}{04.05.08.01}
\label{para:04.05.08.01_StrategischeRessourcenplanung}

% Task
\par Neben dem Sichern der sehr raren fachlichen und technischen internen Ressourcen für \gls{r2h} war die Einbindung externer Dienstleister mit passendem Fachwissen, insbesondere zur \gls{aif}, sowie ausreichender Erfahrung eine zentrale Herausforderung. Meine Aufgabe bestand darin, eine strategische Ressourcenplanung zu entwickeln, um die Projektergebnisse termingerecht und in der erforderlichen Technologie liefern zu können. Dabei musste ich sowohl die internen Kapazitäten als auch die Verfügbarkeit und Leistungsfähigkeit externer Partner berücksichtigen, um eine effektive und effiziente Ressourcennutzung sicherzustellen.

% Action
\par Um die strategische Ressourcenplanung zu entwickeln, führte ich zunächst eine umfassende Bedarfsanalyse durch, um die erforderlichen Fähigkeiten \gls{btp}, \gls{aif} und Kapazitäten für die verschiedenen Projektwellen (vgl. \autoref{tab:wellenzuordnung}) zu identifizieren. Anschließend erstellte ich einen Ressourcenplan, der sowohl interne als auch externe Ressourcen berücksichtigte und auf die Projektmeilensteine sowie die geplanten Wellen abgestimmt war. Dabei legte ich besonderen Wert auf die frühzeitige Einbindung externer Dienstleister, um deren Expertise optimal nutzen zu können und potenzielle Engpässe zu vermeiden. Ergänzend definierte ich klare Priorisierungs- und Steuerungskriterien für den Ressourceneinsatz sowie regelmäßige Überprüfungs- und Anpassungsprozesse, um die Ressourcenplanung flexibel an sich ändernde Anforderungen und Rahmenbedingungen anpassen zu können.

% Result
\par Aufgrund des Paradigmenwechsels von \gls{gfa} mit \gls{aif} zu \gls{bfa} mit \gls{ft} (vgl. \autoref{fig:R2HInterfacesIntegration}) konnten die ursprünglich geplanten externen Ressourcen für die \gls{aif}-Implementierung nicht wie vorgesehen eingebunden werden. Durch die Wiederaufnahme des \gls{f2h}-Projekts als Scopeerweiterung wurde die Ressourcenplanung zwar grundsätzlich verkompliziert; die externen Ressourcen mit \gls{aif}-Know-how konnten jedoch in das separat ressourcengeplante \gls{f2h} (vgl. \autoref{fig:ADNSSTFI}) transferiert werden. Damit konnte vorhandene externe Expertise weiterhin zielgerichtet eingesetzt und die Lieferfähigkeit trotz veränderter technischer Rahmenbedingungen stabilisiert werden.

% Reflect
\par Aus Ressourcenengpässen resultierende Herausforderungen können durch eine strategische Ressourcenplanung zwar nicht vollständig eliminiert, aber durch frühzeitige Bedarfsanalyse, proaktive Einbindung externer Partner und flexible Anpassungsprozesse effektiv gemanagt werden. Grundsätzlich ist es jedoch essenziell, die Ressourcenplanung kontinuierlich zu überwachen und bei Bedarf anzupassen, um auf unvorhergesehene Veränderungen reagieren zu können und die Projektergebnisse dennoch termingerecht und in der erforderlichen Qualität zu liefern.